\section{From the Case V guard to a separating trace bound}
\label{sec:count-v-install-v61}

Here $\ell$ is the original coefficient prime, $O=c.O$, $k$ its residue
field, and $K=\mathbb Q(i)$. The source supplies a specified algebraic closure $\Omega$ of $\mathbb Q$. Write
$\Delta=\operatorname{Gal}(\Omega/\mathbb Q)$.
The algebraic inputs called Lemma M and Lemma M-prime in this section
are the valuation and length results SP1--SP6.

% BEGIN CVI1
\paragraph{CVI1. The Case V hypotheses supply both host predicates.}
\label{claim:CVI1}
Fix a character $\chi$, a host $h$ with $h.\chi=\chi$, a permitted
$t$ with $h.\mathrm{Big}(t)$, an odd prime $\ell$, and
$1\le e<t$. Fix $\psi\in h.X(c,e)$ and suppose that
$h.\mathrm{CaseV}(\ell,e)$ and
$R_1=h.R1(\ell,\psi)\le2^{c_2(\psi)}$.
In particular $e=1$ and $\ell\in B_\chi$. Take $q=h.q$, $o=h.o$,
$N=h.N$; $q$ is the positive conductor of $\chi$, and $o$ its order.
Put $\theta=\operatorname{th1}(\ell,\psi)$ as in HXD1.

Here is the character calculation at this host. Its defining subgroup
is the intersection of the inverse images of $\langle-1\rangle$ at
the dyadic level, $\ker\chi$ at $q$, and the square subgroups at the
two auxiliary odd primes. It is killed by $\psi$. For each odd
$p\mid q$ the full $p$-power in $N$ is the one in $q$, since the
two auxiliary primes are coprime to $q$. The $o$-th power of every
unit supported at $p$ lies in the host subgroup. At each of the two
auxiliary primes outside $q$ the square of every local unit lies in
that subgroup. At the other odd primes outside $q$ there is no
component. Thus the odd $p$-parts of $\psi$, and their prime-to-$\ell$
parts, have orders dividing $2o$. The square of $\theta$ has conductor
dividing $q$: its dyadic component is trivial by construction,
its components at the two auxiliary primes square to one, and the
remaining conductor powers divide $q$. The $\ell$-primary part of
$\psi$ is trivial on the dyadic and auxiliary factors since those
images have order prime to $\ell$. In particular every
$v\in T_b(\ell,\psi)$ divides $\ell q$.

The order calculation of HXD1 uses just the odd prime and the prime-to-$\ell$ order of $\theta$: taking $d$ to be the order of $\ell$ modulo $\operatorname{lcm}(4,\operatorname{ord}\theta)$ gives $n_0=v_2(\ell^d-1)\ge2$. Since $o,q\ge1$ and
$\ell q\ge3$, the original $R_1=16o(\ell q)^2 4^{n_0}$ is at least
$16\cdot9\cdot16>2^{11}$. The guard implies $c_2(\psi)\ge12$
and for every $v\in T_b$,
\[
 8o v^2\le8o(\ell q)^2<R_1\le2^{c_2(\psi)} .
\]
Moreover $16\mid N$ by the original host exponent $\mathrm{ex}\ge2$,
and $\psi(-1)=(-1)^e$ by membership in $h.X(c,e)$. These are the
fields of $\mathrm{HostData}(\ell,e,\psi,o)$, including the
prime-to-$\ell$ component bound above. Since
$2(\ell q)4^{n_0}\le R_1$, the positive integer $q$ and the
conductor divisibility give $\mathrm{XData}(\ell,\psi,q)$.
All calculations allow $\ell=3$ and equality at the $R_1$ boundary.
% END CVI1

% BEGIN CVI2
\paragraph{CVI2. One actual Galois frame and the class-one condition.}
\label{claim:CVI2}
Let $S$ consist of the rational primes dividing $\ell N$, choose
one valuation ring $A_v$ above each, and let $\Gamma=G_S$ be the
quotient of $\Delta$ by the closed normal subgroup generated by all
inertias above primes outside $S$. Use its quotient map $\pi$ and
the actual $D_v=D(A_v)$ and $I_v=I(A_v)$. Write $v_\ell$ for the
place at $\ell$ and use the supplied $S_b(\ell,\psi)$ for the marked
set of bad places. The Galois character $\eta_4$ of $\mathbb Q(i)$
and $\chi_\ell^e\psi_G$ descend to $\Gamma$: at every omitted
prime their finite-level cyclotomic factors are unramified, and their
kernels are closed. Call the descended characters $\eta,\Psi$.
The image of a complex conjugation is the element $c_F$ of the frame.
This gives every field of $\mathrm{IsHostFrame}$. The order of choices
keeps $S,A_v,F$ fixed for the rest of the argument.

The native $F.\mathrm{Ok}$ conditions hold. The residue field is
finite of characteristic $\ell$ by HXD1, and $2$ is a unit. The
quotient map is continuous, and the finite-order character and
all reductions of $\Psi$ modulo powers of the maximal ideal are
locally constant. On $c_F$ their respective values are $-1,1$ by
the parity in CVI1, and $c_F^2=1$. The inertia and decomposition
groups are closed, with inertia normal in decomposition. Their
quotient is generated topologically by residue Frobenius and maps
onto every finite cyclic group with Frobenius mapping to $1$:
the reductions on the finite residue extensions give exactly these
maps. At $T_b=\{v_\ell\}\cup S_b$ the primes are odd (the
dyadic unit group has $2$-power order), so $\eta$ is unramified.
AFE10 provides, in each of these actual $D_v$, an $x$ with
$\eta(x)=1$, $\bar\chi_\ell(x)=1$ and $\bar\Psi(x)^2\ne1$.
Here and below characters on $\Gamma$ are pulled back along $\pi$
when evaluated on $D_v$. Away from $T_b$ the cyclotomic factor
is unramified and the ramified finite part of $\psi$ has order
prime to $\ell$. Prime-to-$\ell$ roots of unity reduce injectively.
Consequently $\Psi=1$ on the prescribed joint inertia kernel
$N_v=I_v\cap\ker\eta\cap\ker\bar\Psi$. Finally $v_\ell\notin S_b$
by definition. These verifications include all sixteen frame conditions.

Both $\Delta$ and $\Gamma$ are profinite. Global finiteness implies
$\Phi_\ell(\Gamma)$. Indeed, for any open $H\le\Gamma$, induce the
trivial $\mathbb F_\ell$-module of $H$ to the finite permutation
module on $\Gamma/H$. Shapiro identifies its continuous $H^1$ with
$\operatorname{Hom}_{\rm cont}(H,\mathbb F_\ell)$; finite coefficients
and $S\supseteq\{\ell\}$ give finiteness of the former group.

For $F.\mathrm{Cl1}$ let $a:\ker\eta\to B$ be any locally constant
homomorphism to an abelian group which kills the specified inertia
conjugates. Pull back to $G_K\subset\Delta$, the inverse image of
$\ker\eta$. The homomorphism has finite image by compactness.
Every inertia of $K$ is an intersection of $G_K$ with an absolute
inertia of $\mathbb Q$. Outside $S$ the latter group dies in
$\Gamma$; inside $S$ it is conjugate to the chosen $I_v$, and the
conjugate kernel condition on $a$ kills the intersection. Thus the
finite abelian extension defined by the pullback is everywhere
unramified. The Gaussian integers are Euclidean; by the unramified
abelian class-field theorem this extension is trivial. Since
$G_K$ surjects onto $\ker\eta$, $a$ is trivial.
% END CVI2

% BEGIN CVI3
\paragraph{CVI3. Local vanishing and the global degree-two restriction.}
\label{claim:CVI3}
For every integer $j$ put $M_j=\mathrm{Gb}(k,j)$, the matrices
supported at entries of weight $j$ in the three-dimensional
residual adjoint action. The dimension is
$m_j=\max(3-|j|,0)$. For $j\ne0$ the AFE10 witness of CVI2 acts
on $M_j$ by $\bar\Psi(x)^j$, so $x-1$ acts invertibly:
only $j=\pm1,\pm2$ can contribute. It acts invertibly on
$M_j^*(1)$ as well, since $\bar\chi_\ell(x)=1$.
Thus at $v\in T_b$, the degree-zero local cohomology of the
module and of its Tate dual vanishes. Local duality gives
$H^2(D_v,M_j)=0$. At $v\in S_b$, the local Euler characteristic
away from $\ell$ then gives $H^1(D_v,M_j)=0$. These are exactly
the two clauses of $F.\mathrm{LocVan}$.

For $j\ne0$ the same argument at $v_\ell$ kills the global
Tate-dual invariants. For $j=0$, choose a Galois automorphism
fixing $i$ and acting nontrivially on $\mu_\ell$. It exists by
the CRT at $4$ and $\ell$ and surjectivity of the cyclotomic
quotient over $\mathbb Q$. On degree-zero matrices its action
is the identity (its $\eta$ is $1$), whereas its action on
$M_0^*(1)$ is a nonidentity scalar. Poitou--Tate therefore makes
\[
             H^2(\Gamma,M_j)\longrightarrow
                     \bigoplus_{v\in S} H^2(D_v,M_j)
\]
surjective for every $j$. Odd residue characteristic makes the
positive-degree archimedean Tate cohomology zero. The local and
global finiteness assumptions apply to the original finite
$\mathbb F_\ell$-modules by restriction of scalars from $k$;
the underlying cochain complexes and their maps are unchanged.
Composition with $\operatorname{Tr}_{k/\mathbb F_\ell}$ identifies the
restricted $k$-linear dual with the $\mathbb F_\ell$-linear dual by
nondegeneracy of the finite-field trace pairing, compatibly with
the cyclotomic twist.
% END CVI3

% BEGIN CVI4
\paragraph{CVI4. The Euler defect for the native framed quotient.}
\label{claim:CVI4}
Write $Z_{L,v}\subseteq Z^1_v$ for the local subspace at weight $j$.
All local coboundaries lie in it: off $v_\ell$ the action on $N_v$
is trivial, and at $v_\ell$ the only additional weight-zero
inertia condition is trivial for a coboundary. Put
$L_v=Z_{L,v}/B^1_v$, $U=\bigoplus_vH^1_v/L_v$.
Localization gives an exact sequence
\[
 0\to H^1_L\to H^1\to U\to
          F.\mathrm{Ob}(j)\to H^2\to\bigoplus_v H^2_v\to0.
\]
Indeed an obstruction representative is a global $2$-cocycle and
a local $1$-cochain at each $v$ whose differential equals its
restriction. Forgetting the local cochains maps onto the kernel of
the last restriction. When the global cocycle is a coboundary,
subtract a global primitive together with its restrictions. What
remains is a tuple of local cocycles modulo $Z_{L,v}$ and global
restrictions, precisely the indicated cokernel. CVI3 supplies
surjectivity at the last term. The numerator defining $F$'s
degree-one quotient is the space of cocycles with precisely these
local conditions, giving $H^1_L$ at the other end. Finite-dimensionality
follows at once. Its Euler difference is
\[
 h^1-h^2+\sum_v(h^2_v-\dim Z^1_v+\dim Z_{L,v}).
\]

For $v\ne v_\ell$ the quotient $I_v/N_v$ has order prime to
$\ell$ by the joint-character argument in HXP1, and $N_v$ acts
trivially on $M_j$. Locally $L_v$ is the kernel of restriction to
$N_v$, equivalently to $I_v$: on $N_v$ coboundaries are zero;
if a class vanishes there, its restriction to $I_v$ inflates from a
finite group of invertible order and vanishes by averaging. The
unramified local space has dimension $h^0_v$, since the unramified
quotient is procyclic and its first cohomology on $M_j^{I_v}$
is the Frobenius cokernel, of the same dimension as the kernel.
It follows that $\dim Z_{L,v}=m_j$, cancelling the local term by
the Euler formula $h^1_v=h^0_v+h^2_v$.

Let $a_j=\dim M_j^{c_F}$. Negative transpose followed by reversal
pairs the entries $(a,b)$ and $(2-b,2-a)$; its fixed space has
dimension $a_j=1$ for $|j|\le1$ and $0$ otherwise, since the
characteristic is odd. The global Euler formula gives
$h^1-h^2=h^0+m_j-a_j$ and the $\ell$-adic local formula gives
$h^1_\ell=h^0_\ell+h^2_\ell+m_j$. For $j<0$ the local
condition is $B^1_\ell$; both global and distinguished local
invariants vanish by CVI3, so the difference is $-a_j$. For
$j>0$ the local condition is all of $Z^1_\ell$, $h^2_\ell=0$,
and the difference is $m_j-a_j$. For $j=0$ it is unramified,
with $\dim Z_{L,\ell}=m_0$ by the same Frobenius computation
($I_\ell$ acts trivially). The global invariants at weight zero
are exactly $M_0^{c_F}$ since the sole action on the diagonal
comes from $\eta$. The difference is zero. Thus the integer
$\dim H^1_L-\dim F.\mathrm{Ob}(j)$ equals $-1$ at $j=-1$,
$1$ at $j=1,2$, and zero otherwise. This is the original $F.\mathrm{Count}$.
% END CVI4

% BEGIN CVI5
\paragraph{CVI5. The finite list obtained from a compatible weighted model.}
\label{claim:CVI5}
Fix the original $\iota:\overline{\mathbb Q}_\ell\simeq\mathbb C$.
From $\mathrm{ThmS6\_s}$ take its $g_1\in\Gamma$, marked on
the frame of CVI2. Suppose the source's native marked universal
pair $(R,u)$ and its weighted deformation comparisons described below
are supplied. Then the two algebraic lemmas give an integer $V$ and
a finite list $r_i$ of marked characteristic-zero points of type I
satisfying $B'_s(K,\iota)$ on $\Delta$ and pairwise nonisomorphic
on $\ker\eta$, with
\[
  \mathrm{shaLen}_F(b)\le V\quad(0\ne b\in O),\qquad
  V\le2 e_O\sum_i v_\ell\bigl(
  \operatorname{tr}r_i(\phi)-\operatorname{tr}r_z(\phi)\bigr).
\]
Here $e_O=v_O(\ell)$, $v_\ell(\ell)=1$,
$r_z$ denotes the original diagonal base representation, and
$\phi$ is the single separating element supplied by
$\mathrm{HostNumber\_s}$ for $g_1$. We specify the comparisons used
in this implication to identify the exact scope of the model assumption.

Take an adapted minimal weighted power-series presentation $W\to R$
with origin $u\mapsto r_z$. It must identify its weight parts with
the native parts of the universal representation. The variables
are marked matrix-entry differences from the base representation,
of weights in $[-2,2]$. Write $d_j$ for their
numbers. The native obstruction-to-relations and tangent-to-generators
comparisons give homogeneous kernel generators, at most
$b_j=\dim F.\mathrm{Ob}(j)$ in weight $j$, and
$\dim H^1_L(j)\le d_j$. CVI4 then gives
$b_j+\mathrm{defect}(j)\le d_j$. The nonpositive tangent
comparison takes $F.\mathrm{VanE}$ to the assertion that an
$O$-valued coordinate vector of any weight $j\le0$ killing the
linear parts of these relations is zero. These are comparisons
of the actual framed representations with the chosen coordinates.

The needed point comparisons are as follows. On the graded
entry ring $\mathrm{Sring}(u)$ each geometric point $\xi$ gives
the representation $r_\xi$ by evaluating all entries, with its
original right component. Finiteness of the weight-zero module
gives a marked characteristic-zero point after permitted torus
rescaling. Vanishing on all negative or all positive weights
corresponds to the two triangular flags. Vanishing on both gives
the origin, by the native diagonal-point comparison using $F.\mathrm{Cl1}$.
A point having a nonzero coordinate on each side gives type I,
by the native two-sign comparison using $F.\mathrm{LocVan}$
and $F.\mathrm{Xpp}$. These comparisons concern the
characteristic-zero points of this ring, including points that
require rescaling to an integral representation. On these points,
isomorphism over $\ker\eta$ is equivalent to the prescribed
diagonal torus orbit, and the native tangent comparison sends
rigidity to the formula $D(z)=c\,j\,\xi(z)$ on a homogeneous
element of weight $j$ for every point derivation $D$.

Apply $\mathrm{ThmS6\_s}$ to this universal pair. It gives the
finiteness just used and $B'_s(K,\iota)$ for every such type-I
point. Here $B'_s$ means all three exact requirements of
\texttt{BPrime\_s}, including irreducibility on open subgroups
of the Gaussian kernel, irreducibility on the intersection
with the full cyclotomic kernel, and $\mathrm{AutOnAlt}$ with
sign $-1$ for the pulled-back matrix representation. Hence
$\mathrm{HostRigid\_s}$ applies to each point with these same
parameters. Every two-sided point is separated by the
given $\phi$. The degree-zero element testing this is exactly
\[
 a=\operatorname{tr}u(\phi)-\operatorname{tr}r_z(\phi).
\]
Each diagonal entry has weight zero, the origin evaluates this
difference to zero, and $\xi(a)$ is its stated trace difference.

For precision about the algebraic step, if $d_{-1}=0$, the
injection of the next paragraph has zero image for every $b$,
so take $V=0$ and the empty list. Otherwise pad the homogeneous
relations by zero to exactly $d_{-1}+1,d_{-2},d_1-1,d_2-1,d_0$ equations of the respective weights $-1,-2,1,2,0$, as permitted by the five inequalities.
Nonpositive tangent vanishing means the respective linear forms
span the cotangent parts after inverting a uniformizer, by finite
dimensional linear duality. The weight-zero quotient is finite by
the marked theorem. The diagonal point comparison forces a fixed
point killing both signs to be the origin. The rigidity formula makes
two derivations at a two-sided point linearly dependent, since each
is a scalar multiple of the Euler derivation. The trace verifies the
last nonvanishing hypothesis. In making this identification one uses
the free weighted graded-algebra presentation with its kernel cut
out by the specified homogeneous equations and the original origin.
SP6, with valuation $e_O v_\ell$, bounds its finite weight-minus-one
cotangent length $V$ by $2 e_O\sum_i v_\ell(\xi_i(a))$. One may select
one representative per orbit from the finite cover and apply SP6
to the thinned cover; the estimate holds for every cover. The above
orbit comparison yields the required matrix nonisomorphism.

Finally use the native integral coordinate comparison: for each
$0\ne b\in O$ there is an additive injection
$F.\mathrm{ShaO}(b)\to(O/(b))^{W.\sigma}$ supported on weight
$-1$, annihilating every linearized defining relation. Choosing
arbitrary integral lifts of its coordinate values gives derivations
of the free weighted algebra at the origin; their reductions kill
the weight-minus-one equations. SP1 therefore bounds the length
of this module by the same $V$. The factor $e_O$ enters the upper
valuation bound exactly once. This derives the two displayed estimates
from the stated native model comparisons.
% END CVI5

% BEGIN CVI6
\paragraph{CVI6. Uniform exponent and the literal CountV bound.}
\label{claim:CVI6}
Assume the host inputs in the strengthened source signature and the
native marked-model comparisons of CVI5 for the frames of CVI2.
The resulting estimate is the full original $\mathrm{CountV}$.
Indeed choose the Gaussian field first. $\mathrm{HostNumber\_s}$
then gives one $C\in\mathbb N$ before any $\Omega,\ell,c,N,e,\psi$,
frame, coefficient embedding or marked element. For each of the
later instances choose $\iota$ and take the frame of CVI2, its
$g_1$ from $\mathrm{ThmS6\_s}$, and the $\phi$ of CVI5. This $\phi$
is fixed before every finite collection of points. The same
$\mathrm{HostNumber\_s}$ upper bound therefore applies to the
particular irredundant list of CVI5. With $X=\ell N e$ it says
\[
 \Bigl(\sum_i v_\ell(\operatorname{tr}r_i(\phi)-
                  \operatorname{tr}r_z(\phi))\Bigr)\log\ell\le X^C .
\]
VLT3 gives $\mathrm{Vlen}\le V<\infty$ at these original
frame lengths. Its normalization is division by the positive
integer $e_O$, so
\[
 \mathrm{Vn}\log\ell\le(V/e_O)\log\ell
 \le 2\Bigl(\sum_i v_\ell(\operatorname{tr}r_i(\phi)-
                         \operatorname{tr}r_z(\phi))\Bigr)\log\ell
 \le 2X^C\le X^{C+1}.
\]
The last inequality uses $X\ge2$, which holds for every permitted
prime, positive level and positive exponent. The same $C+1$ works
for every fixed target character $\chi$, including every $h$ with
$h.\chi=\chi$ and every $t$ in its original scope. CVI1 uses the
weak side $R_1\le2^{c_2(\psi)}$; the separate
$\mathrm{CountVs}$ concerns the strict reverse inequality.
% END CVI6

\paragraph{Boundary of this reconstruction.}
The strengthened formal implication takes exactly
$\mathrm{ThmS6\_s}$, the two algebraic lemmas, local finiteness,
local Tate duality, local Euler characteristic, global finiteness,
global Euler characteristic, normalized Poitou--Tate,
$\mathrm{HostVanE}$, $\mathrm{HostXpp}$,
$\mathrm{HostRigid\_s}$ and $\mathrm{HostNumber\_s}$.
The selected graph separately tracks the frame-length work;
the formal implication constructs its requisite native deformation
objects internally. CVI1--4 establish arithmetic, class-group and
cohomological components; VLT1--3 give the actual length comparison.
CVI5--6 explicitly retain the native marked universal pair,
its adapted presentation and equations, the point, tangent and
orbit comparisons, and the integral Selmer coordinate injection.
The pinned source supplies these in the coordinate proof of
\texttt{abstract\_count}. The readable specialization of the
previous general and one-character presentations to all these
native comparison statements is still outstanding here.
